\documentclass[11pt]{article}
\usepackage[utf8]{inputenc}
\usepackage[T1]{fontenc}
\usepackage{lmodern}
\usepackage[margin=1in]{geometry}
\usepackage{amsmath,amssymb}
\usepackage{textcomp}
\usepackage[dvipsnames]{xcolor}
\usepackage{enumitem}
\usepackage{parskip}
\usepackage{titlesec}
\usepackage[colorlinks=true,linkcolor=blue,citecolor=blue,urlcolor=blue]{hyperref}

\titleformat{\section}{\large\bfseries}{\thesection}{1em}{}
\titleformat{\subsection}{\normalsize\bfseries}{\thesubsection}{1em}{}
\titleformat{\subsubsection}{\normalsize\bfseries\itshape}{}{0em}{}

\title{\textbf{Where Thinking Starts to Hurt:\\ Measuring Difficulty-Conditioned Tipping Points in Test-Time Reasoning}}
\author{}
\date{}

\newcommand{\ph}{\textit{[placeholder --- fill in after experiments are run]}}

\begin{document}
\maketitle
\vspace{-2em}

\begin{quote}
\small\textit{Assembled from the six confirmed stage outputs of an \texttt{auto-research-discussion} run (seed topic B4.3), post research-taste audit (\textcolor{ForestGreen}{\textbf{PASS}}). Status: \textbf{specification / paper draft --- no experiments have been run.} All numbers below are pre-registered design targets, not results.}
\end{quote}

\begin{abstract}
Reasoning models expose ``thinking budget'' as a controllable resource: spend more test-time compute, get more accuracy --- or so the assumption goes. This assumption hides an asymmetry that recent evidence has begun to expose: additional thinking tokens do not merely stop helping past some point; the same generative process that produced a correct answer can overwrite it, so the marginal return can turn negative rather than flatten, and this failure mode is difficulty-conditioned rather than uniform. Inverse-scaling evaluation suites construct settings where longer reasoning measurably deteriorates accuracy, and a systematic budget sweep now reports diminishing marginal utility, correct-answer abandonment, and an optimal thinking length that varies with difficulty. What none of this evidence --- nor the dense adjacent literature of stopping rules, pre-hoc difficulty probes, and length-shortening training objectives built on the same presupposition --- publishes is the boundary itself: a statistically localized, per-difficulty-stratum tipping point with stated uncertainty, validated against manipulation artifacts, and tested for cheap pre-hoc predictability. We build a measurement harness that triangulates an operational difficulty proxy against two independent stratifications, localizes a per-stratum tipping point with a paired-bootstrap change-point test replicated across three independent length-control manipulations, and evaluates whether cheap signals --- output-entropy trajectories and early-answer consistency, computable from a short prefix of the reasoning trace --- predict that tipping point beyond a coarser difficulty-bin decision and a length-only null. We report whether, and where, the marginal return of test-time reasoning flips sign as a function of difficulty, and how much of the achievable accuracy-cost tradeoff a cheap pre-hoc predictor of that boundary actually recovers.
\end{abstract}

\section{Introduction}

Test-time ``thinking'' has become one of the few knobs a reasoning model exposes after training is finished: ask it to generate a longer chain of thought before answering, and accuracy on hard problems often rises. This is cheap to implement --- no retraining required --- and it is tempting to treat as a free lunch, which is exactly why it is worth checking whether it actually is one.

Yet a central assumption behind budget-allocation decisions for these systems has only recently begun to crack, and only at the level of existence evidence: that the marginal benefit of more thinking merely \emph{flattens} once it stops helping, rather than turning negative. If instead the same reasoning process that produced a correct answer can, given enough additional tokens, talk itself out of that answer, then a policy that spends a fixed or ever-growing thinking budget is not simply wasteful on the problems that don't need it --- it is actively destructive on exactly those problems. Where this happens, as a function of how difficult a problem is --- with stated uncertainty, robustly to how thinking length is manipulated, and predictably in advance --- is the question this paper measures rather than assumes.

The nearest existing work approaches this question from adjacent but incomplete angles, and it is worth naming the deltas plainly rather than surveying broadly first. A theoretical account of chain-of-thought length \cite{WhenMoreIsLess6QDF} derives an inverted-U relationship between reasoning length, difficulty, and accuracy, predicting exactly the kind of turn this paper looks for --- but validates that prediction only on synthetic tasks where difficulty is a variable the experimenters construct, not a quantity measured on natural benchmarks. An empirical study of ``overthinking'' \cite{Chen2024Overthinking} documents that reasoning models expend redundant solution rounds on easy problems with negligible or negative benefit, and an inverse-scaling evaluation suite \cite{InverseScalingTTC2507} constructs task categories where longer reasoning deteriorates accuracy outright --- direct evidence the phenomenon exists, in natural and constructed settings alike. Closest of all, a recent systematic study of test-time compute scaling \cite{WhenMoreThinkingHurts2604} reports that marginal returns diminish substantially at higher budgets, that extended reasoning is associated with abandoning previously correct answers, and that the optimal thinking length varies with problem difficulty. This paper concedes all of that as prior evidence; what even this closest neighbor stops short of is the object this paper is built to produce: a per-difficulty-stratum boundary that is statistically localized with stated uncertainty (or explicitly right-censored), shown to replicate across independent length-control manipulations rather than being an artifact of one, and tested for predictability from signals available before the full budget is spent. Between a theoretical shape validated only on synthetic difficulty, existence evidence without a localized boundary, and a difficulty-varying optimum reported without uncertainty, manipulation control, or a pre-hoc predictor, the gap this paper closes is the boundary's measurement, not its existence.

Turning this gap into something measurable surfaces its own difficulties. Difficulty itself is not directly observable for a fresh problem instance --- any stratification needs an operational proxy, and the obvious one (a model's own solve rate on a small sample) risks encoding that model's idiosyncrasies rather than the task's intrinsic difficulty. Locating exactly where a noisy accuracy curve's slope turns negative is a harder statistical target than estimating accuracy itself, since it depends on a second-order feature of already-noisy data. The two ingredients that would make a cheap pre-hoc signal useful --- that it predicts the boundary, and that it does so with information beyond what a cheaper, coarser difficulty classifier already provides --- are each easy to claim and hard to actually demonstrate, since entropy and consistency signals mechanically co-vary with how much text has already been generated. And because an externally truncated or artificially extended reasoning trace is not obviously the same object as a trace that stopped on its own, any claim that ``more thinking hurts'' risks being a manipulation artifact rather than a property of the reasoning process itself --- a confound that, as this paper's own design discipline makes explicit, admits mitigation but not full elimination by any experiment run here.

Our approach turns each of these obstacles into an explicit design commitment rather than a hedge. Difficulty is triangulated against two independent stratifications before any curve is trusted, converting model-idiosyncrasy from an assumed risk into a measured and reportable one. The tipping point is localized with a change-point test paired at the level of individual items, giving a second-order statistic the statistical power a raw curve comparison would lack, and is only reported as confirmed where at least two of three independent length-control manipulations agree in sign --- an explicit, disclosed mitigation of the intervention-versus-emergence confound, not a claim to have eliminated it. A cheap-signal predictor of the tipping point is evaluated specifically against two decoupling controls --- a coarser difficulty-bin classifier and a length-only null --- so that any claimed information content survives the most obvious alternative explanations rather than merely beating a strawman.

We evaluate this design on natural math and science benchmarks spanning an easy-to-hard range, alongside a synthetic tunable-difficulty family that guarantees at least some strata where a tipping point is reachable within an affordable length budget, comparing against fixed-budget, uniform-budget-forcing, coarse-probe, and oracle-stop alternatives at matched cost.

In summary, we contribute:
\begin{itemize}[leftmargin=1.4em]
\item A difficulty-triangulated measurement design that treats a task's difficulty stratum as a validated, cross-checked quantity rather than an assumed one, reporting agreement (or disagreement) across independent proxies rather than trusting a single model's self-assessment.
\item A paired, item-level change-point test that statistically localizes --- or explicitly right-censors --- the length at which the marginal accuracy return of additional thinking turns negative, replicated across three independent length-control manipulations to separate a genuine property of the reasoning process from a manipulation artifact.
\item A cheap, pre-hoc predictor of the tipping point, built from signals computable within a short trace prefix, evaluated against a coarser difficulty-bin control and a length-only null to establish whether its information is non-redundant rather than merely correlated with easier-to-obtain signals.
\item An empirical account of whether, and how, the tipping point shifts with task difficulty across natural and synthetic domains, and how much of the achievable accuracy-cost tradeoff a cheap pre-hoc predictor of that boundary actually recovers relative to fixed budgets and a confidence-based early-exit baseline.
\end{itemize}

\section{Problem Definition}

\noindent\textbf{Formal setup.} Consider a reasoning model $f$ with an externally controllable thinking budget, and an item $x$ with gold label $y$ and difficulty stratum $s(x) \in \{1,\dots,5\}$. A rollout at truncation length $L$, drawn from a fixed grid $L \in \mathcal{L} = \{256, 512, 1\text{k}, 2\text{k}, 4\text{k}, 8\text{k}, 16\text{k}\}$ thinking tokens, yields correctness $Z(x,L,r) \in \{0,1\}$ for rollout $r$. Stratum accuracy is $A_s(L)$, the mean of $Z(x,L,r)$ over items in stratum $s$ and rollouts $r$. The object of study is, per stratum, whether there exists a length $L_k \in \mathcal{L}$ at which the forward difference $A_s(L_{k+1}) - A_s(L_k)$ becomes significantly negative --- the tipping point $L^*_s$, reported as the smallest such $L_k$, or as right-censored if no such $L_k$ exists within $\mathcal{L}$ --- and whether $L^*_s$ can be predicted, before it is directly observed, from features computable within a short prefix of the maximum length in $\mathcal{L}$.

\noindent\textbf{Why this is hard.} Difficulty itself is not an observable property of $x$; $s(x)$ must be assigned by a computable proxy, and any such proxy risks encoding the assigning model's own idiosyncrasies rather than the item's intrinsic difficulty, so a stratification must be validated against independent evidence before curves computed over it can be trusted. $A_s(L)$ is a Bernoulli mean over stochastic rollouts, so detecting where its \emph{forward difference} changes sign is a second-order statistical target that requires substantially more samples per cell than estimating $A_s(L)$ itself would. Two natural predictive signals for $L^*_s$ --- entropy over a trace's early tokens, and agreement among early truncated-answer probes --- are each entangled with how much text has been generated so far, so demonstrating that either carries information about $L^*_s$ requires ruling out that the information is simply a repackaging of stratum membership or trace position. Finally, because $L$ is imposed externally (by truncation or forced continuation) rather than emerging from the model's own unconstrained behavior, a truncated-and-forced trace and a naturally-stopped trace are not obviously the same object, so a finding obtained under one length-control manipulation alone cannot, by itself, rule out that the finding is an artifact of that manipulation rather than a property of $L^*_s$.

\section{Method Overview}

At a conceptual level, this is a single measurement pipeline with five cooperating parts, each addressing one obstacle named above. A difficulty-triangulation step turns an unvalidated proxy into a cross-checked stratification before any curve is computed. A paired change-point test turns per-item rollout data, collected under a controlled length grid, into a statistically localized (or explicitly right-censored) tipping point per stratum. A cross-arm agreement requirement turns three independently-manipulated sets of length-controlled traces into a single, disclosed-confidence finding about whether a detected tipping point reflects the reasoning process rather than one manipulation's artifact. A decoupled cheap-signal predictor turns a short prefix of a trace into a predicted tipping point, evaluated specifically against two controls designed to rule out that its apparent information is actually difficulty-bin membership or trace-position in disguise. And a trace-content mechanism probe turns the localized tipping point from a where-claim into a why-claim, checking whether the accuracy drop past the boundary coincides with a rise in self-overwriting (correct-to-incorrect) answer flips and verification/backtracking activity, rather than a diffuse, mechanism-free decay.

The \textbf{difficulty triangulation} step computes a low-budget solve-rate proxy per item, quintile-bins it, and requires rank-correlation agreement with two independent stratifications (dataset-native tags, a second model's own solve rate) before a stratum is treated as difficulty-relevant. The \textbf{paired change-point test} takes per-item rollout outcomes at each grid length and localizes the smallest length with a significant negative forward difference, using an item-paired bootstrap rather than an aggregate-curve comparison. The \textbf{triangulated intervention design} runs the same items under three length-control arms --- truncation, forced continuation, and natural stopping --- and reports a tipping point as confirmed only where at least two arms agree in sign. The \textbf{decoupled cheap-signal predictor} computes entropy-trajectory and early-consistency features from a short trace prefix and predicts the tipping point, benchmarked against a reimplemented coarse difficulty-bin classifier and a position-controlled length-only null to isolate genuine trajectory-level information from cheaper or spurious alternatives.

\section{Methodology}

\subsection{Difficulty Triangulation}
\textbf{Challenge.} Difficulty is not directly observable for a fresh problem instance, and the most natural proxy --- a model's own solve rate on a small sample --- is defined with respect to that same model's later behavior, risking a stratification that reflects one model's idiosyncrasies rather than the task's intrinsic difficulty.

\textbf{Insight.} The model-relativity risk cannot be eliminated analytically, since there is no proxy-free ground truth for difficulty to check against --- but it can be \emph{measured} rather than assumed away, by requiring agreement across independently-sourced stratifications instead of trusting any single one.

\textbf{Method.} Each item's difficulty proxy is its solve rate over eight rollouts truncated at 512 thinking tokens; items are quintile-binned on this score. Before any downstream curve is computed, the resulting stratification's rank correlation is checked against two independent stratifications --- the benchmark's own native difficulty tags where available, and a second model's own low-budget solve rate --- with a pre-specified agreement threshold. Strata that fail both checks are reported as unreliable and excluded from downstream tipping-point claims rather than silently retained.

\subsection{Paired Change-Point Test}
\textbf{Challenge.} A single stratum's accuracy at a given length is a noisy Bernoulli mean; detecting where its \emph{slope} turns negative is a second-order statistical target that ordinary curve-fitting or eyeballing cannot reliably localize, and an unlocalized ``it looks like it eventually goes down'' claim is not falsifiable.

\textbf{Insight.} Pairing each bootstrap resample at the level of individual items --- rather than resampling the aggregate curve --- gives the slope-sign test power to detect a genuine forward-difference change without smoothing over exactly the effect being sought.

\textbf{Method.} At each grid length, rollouts are collected per item under a fixed harness. Per stratum, an item-paired bootstrap (10,000 resamples, one-sided, Holm-corrected across strata and length-control arms) tests whether the forward difference between adjacent grid lengths is significantly negative. The tipping point is reported as the smallest grid length with a significant negative step; strata with no such length within the studied grid are reported as right-censored rather than assigned a value by extrapolation. A trend test across ordered difficulty strata (with censored strata assigned maximum rank) additionally checks whether the tipping point itself shifts systematically with difficulty.

\textbf{Placebo check.} Before any stratum's tipping point is trusted, the identical bootstrap and change-point procedure is run on a null-by-construction control: grid-length labels are randomly permuted within each item's rollout set (severing any true length-accuracy relationship while preserving the marginal accuracy distribution), and the same Holm-corrected test is applied. A non-null detection under this permutation --- a ``significant'' flip found where none can exist by construction --- would indicate a pipeline artifact (e.g.\ multiple-comparison leakage or a miscalibrated bootstrap) rather than a property of the reasoning process, and would block scoring of every downstream hypothesis until fixed, mirroring the placebo-first discipline applied to the length-control arms themselves.

\subsection{Triangulated Intervention Design}
\textbf{Challenge.} Externally imposing a thinking length --- by truncating a trace or forcing it to continue --- is not obviously the same intervention as a model naturally arriving at a shorter or longer trace on its own, so a negative-slope finding obtained under a single length-control manipulation risks being an artifact of how that manipulation was applied rather than a property of the reasoning process itself.

\textbf{Insight.} No single manipulation realizes the exact counterfactual of the same, unmodified model choosing to think longer on the same item --- full causal identification is not achievable by any design available to a measurement-only study. What is achievable, and what this design targets, is evidence that a finding survives across manipulations that differ in exactly how they impose length, which is a disclosed mitigation of the confound, not a claim to have eliminated it.

\textbf{Method.} Three arms are run on identical items: truncation (the thinking segment is cut at a fixed length and an answer is forced), forced continuation (the natural end-of-thinking signal is suppressed and the trace is extended, in the style of prior budget-forcing work), and natural stopping (unconstrained generation, logged observationally). A tipping point is reported as confirmed only where at least two of the three arms agree in sign; detections that appear in only one arm are reported separately as unconfirmed rather than folded into the headline finding.

\subsection{Decoupled Cheap-Signal Predictor}
\textbf{Challenge.} Signals available early in a trace --- output-token entropy, agreement among early truncated answers --- are indirect proxies for eventual correctness whose apparent predictive value could simply be a repackaging of which difficulty bin an item falls in, or of how much text has been generated so far, rather than genuine trajectory-level information about the tipping point.

\textbf{Insight.} A predictor's information is only demonstrably non-redundant with cheaper alternatives if it is evaluated specifically against those alternatives under matched conditions --- beating an uninformative comparison proves nothing about whether the trajectory itself carries anything a coarser or positional signal doesn't already give away.

\textbf{Method.} From the first 256 thinking tokens plus four truncated-answer probes, entropy-trajectory statistics (mean, slope, and windowed maximum of token-level predictive entropy) and an early-consistency score (agreement rate among the four probe answers) are computed and fed to a lightweight predictor of the tipping point, trained on one split of strata and evaluated on held-out strata and held-out task domains. Its predictive value is assessed against two decoupling controls under matched holdouts: a reimplemented three-class difficulty-bin classifier from prefill hidden states (the coarse-bin control), and a position-controlled length-only null feature set (the trace-position control). Ground-truth correctness is never required at prediction time. Both hypotheses evaluated under this design are scoped to held-out strata and domains within the models tested; cross-model-family transfer of the predictor is an explicitly declared limitation, not a silent omission.

\subsection{Trace-Content Mechanism Probe}
\textbf{Challenge.} Localizing \emph{where} the marginal return turns negative (the tipping point itself) is a claim about existence and location, not about cause --- without a mechanism check, ``more thinking hurts'' could reflect any number of unrelated processes (e.g.\ context-window degradation, repetition collapse) rather than the specific self-overwriting story motivating this paper (Introduction): that the same generative process which reached a correct answer can talk itself out of it.

\textbf{Insight.} A mechanism claim is only earned by showing that the answer-flip behavior itself changes shape at the localized boundary, not merely that aggregate accuracy does --- a rise in correct-to-incorrect flips co-occurring with elevated verification/backtracking activity right at the tipping point is a qualitatively different, and falsifiable, signature from a generic decay.

\textbf{Method.} Restricted to strata with a detected tipping point $L^*_s$ (Methodology, Paired Change-Point Test), compare reasoning-trace content just before vs.\ just after $L^*_s$ across rollouts along three measures: (i) the self-correction rate --- the fraction of rollouts where the truncated-answer probe's answer changes between adjacent grid lengths, split by whether the change is correct$\to$incorrect or incorrect$\to$correct; (ii) verification/backtracking step density per 256-token window (a keyword-and-structure heuristic, e.g.\ ``wait,'' ``let me check,'' re-derivation), before vs.\ after $L^*_s$; (iii) the correlation between the entropy-trajectory slope feature (Decoupled Cheap-Signal Predictor) and the sign of the answer flip at $L^*_s$. The pre-registered prediction is that the correct$\to$incorrect flip rate rises sharply past $L^*_s$ relative to before it, co-occurring with elevated verification-step density, rather than a diffuse, mechanism-free accuracy decay.

\section{Experiments}

\subsection{Experimental Setup}

\textbf{Datasets.} MATH-500 (native Level 1--5 difficulty tags, auto-gradable numeric answers) is the pilot-stage domain, chosen because closed-form grading keeps the triangulation and change-point machinery cheap to validate first. GSM8K serves as an easy-stratum anchor, checking that the lowest-difficulty end of the curve is not an artifact of MATH-500's own level-1 items being unrepresentatively hard. AIME 2024/2025 serve as a hard-stratum anchor, used for curves rather than per-item statistics given their small annual item count, to probe whether the hardest strata are right-censored within the studied grid. GPQA-Diamond, a non-mathematical multiple-choice domain, stress-tests whether the predictor's held-out-domain claims hold outside math-heavy, free-response auto-grading, rather than being observed only within the math corpora above. A synthetic tunable-difficulty family (e.g., Countdown or multi-digit arithmetic with a controlled operand count) guarantees at least some strata where the tipping point is reachable at low budget, de-risking the possibility that every natural-benchmark stratum's curve is still rising at the largest studied length.

\textbf{Baselines} (sanity-check comparators giving the recovery and informativeness claims a scale, not competitors to beat --- this is a measurement study, and no research question here is framed as ``our procedure beats baseline X''). Fixed thinking budgets at each grid length, including an always-max-budget policy, form the no-stratification control: if the best fixed budget already sits within one percentage point of oracle-stop at matched cost, the boundary map has no allocation value, which is itself an informative outcome rather than a failed comparison. Uniform budget forcing, applied identically across all difficulty strata regardless of stratum, gives the difficulty-conditioning claim a concrete ``one global policy'' alternative to rule out. A reimplemented DEER-style confidence-exit policy \cite{DEERNpU7} --- stopping when a trial answer's confidence crosses a threshold, with no difficulty stratification --- is the within-trace dynamic-stopping comparator: it tests whether an existing, already-published training-free stopping rule already captures the tipping point implicitly, without ever localizing or reporting it. A reimplemented three-class prefill-difficulty probe and a position-controlled length-only null feature set are the two decoupling controls for the predictor's non-redundancy claim. Per-stratum oracle-stop, computed in hindsight, establishes the ceiling against which ``substantial fraction recovered'' is defined.

\textbf{Compute budget (design envelope).} The pilot is bounded above by 500 items $\times$ 7 grid lengths $\times$ 8 rollouts $\approx$ 28k rollouts per interventional arm; since truncation caps thinking at each grid length, one full pass over the grid costs at most the grid sum ($\approx$ 32.5k thinking tokens) per item-rollout pair, giving $\approx$ 0.13B thinking tokens per arm and $\approx$ 0.26B across the two interventional arms (natural stopping is logged from the same generations). The scale-up ($\geq$3 models $\times$ $\sim$2{,}500 items $\times$ 7 lengths $\times$ 16 rollouts $\times$ $\geq$2 arms) is bounded by $\approx$ 5--10B thinking tokens. Both figures are declared here as the envelope this pre-registration commits to --- days-to-weeks on a self-hosted multi-GPU node for open-weight models, or a mid-four-figure dollar API budget --- so that feasibility is auditable before any result exists.

\textbf{Metrics.} Stratum accuracy at each length and its bootstrap-estimated forward difference; the localized tipping point (or right-censoring indicator) with a bootstrap confidence interval; per-feature AUROC for classifying the sign of marginal gain beyond the probe point; oracle-gap recovery, defined as the accuracy gain of predicted-stop over the best fixed budget divided by the accuracy gain of oracle-stop over the best fixed budget, at matched mean token cost; and cross-stratification rank-correlation agreement for the difficulty proxy.

\textbf{Research questions:}
\begin{enumerate}[leftmargin=1.6em]
\item RQ1 --- Does the difficulty-proxy stratification agree with two independent stratifications?
\item RQ2 --- Does the marginal accuracy return of thinking turn significantly negative in some difficulty strata?
\item RQ3 --- Does the tipping point shift systematically with difficulty?
\item RQ4 --- Do cheap trajectory-level signals carry information about the tipping point beyond a length-only null?
\item RQ5 --- Does the trajectory-level predictor's information about the tipping point exceed what a coarser difficulty-bin decision already recovers?
\item RQ6 --- Does a detected tipping point replicate across independent length-control manipulations?
\item RQ7 --- Does the tipping point coincide with a mechanistic signature of self-overwriting, or only with a generic accuracy decay?
\end{enumerate}

\subsubsection{RQ1: Does the difficulty-proxy stratification agree with two independent stratifications?}
\textbf{Setup.} For every tested (model $\times$ dataset) cell, compute Spearman rank correlation between the primary model's low-budget-solve-rate quintiles and (a) the dataset's native difficulty tags where available, and (b) a second model's own low-budget solve-rate quintiles.

\textbf{Expected outcome.} Both correlations reach $\rho \geq 0.4$; strata that fail both checks are reported as unreliable and excluded from downstream tipping-point claims rather than silently retained.

\textbf{Falsification condition.} Refuted if either correlation falls below 0.4 in a majority of tested cells --- that would mean the difficulty proxy is closer to a model-specific artifact than to a task-intrinsic quantity, undermining every downstream stratum-conditioned claim built on it.

\textbf{Hypothesis (plain language).} The difficulty proxy is a validated, cross-checked quantity, not an unexamined artifact of the one model used to compute it.

\textbf{Metrics.} Spearman $\rho$ against native tags; Spearman $\rho$ against a second model's solve rate; fraction of strata excluded as unreliable.

\textbf{Result.} \ph

\subsubsection{RQ2: Does the marginal accuracy return of thinking turn significantly negative in some difficulty strata?}
\textbf{Setup.} Pilot stage: one model, MATH-500, 7 grid lengths, 8 rollouts per (item, length). Scale-up: at least three models spanning scale, four datasets, 16 rollouts per cell. Item-paired bootstrap (10,000 resamples, one-sided, Holm-corrected across 5 strata $\times$ 3 arms) on the forward accuracy difference between adjacent grid lengths.

\textbf{Sample-size rationale.} Pilot cells are sized to detect a 5-percentage-point accuracy drop between adjacent grid lengths at $\alpha = 0.05$, power 0.8 --- a pilot-only target, deliberately coarser than the confirmatory bar. Because the expected outcome below requires a $\geq$3-percentage-point peak-to-max contrast under Holm correction across 5 strata $\times$ 3 arms, scale-up cells are sized for that 3-point contrast at the corrected $\alpha$ --- this, not convenience, drives the move from 8 to 16 rollouts per cell --- and an underpowered pilot result triggers the pre-committed enlargement rather than a null claim.

\textbf{Expected outcome.} In at least 2 of 5 difficulty quintiles, the peak-to-max accuracy contrast is significantly negative ($\geq$3 percentage points, Holm-corrected $p < 0.05$).

\textbf{Falsification condition.} Refuted if no quintile, dataset, or model shows a significant negative forward difference anywhere in the studied grid under any length-control arm --- in which case the paper's own headline reframes from a boundary map to a well-powered null bounding where such a boundary cannot lie within practical budgets, which is itself the reportable finding.

\textbf{Hypothesis (plain language).} For at least some difficulty levels, additional thinking tokens do not merely stop helping --- they measurably hurt.

\textbf{Metrics.} Per-cell accuracy and its bootstrap confidence interval; peak-to-max contrast and its significance; count of quintiles with a detected flip.

\textbf{Result.} \ph

\subsubsection{RQ3: Does the tipping point shift systematically with difficulty?}
\textbf{Setup.} Among strata with a detected tipping point (RQ2), a Jonckheere-Terpstra trend test on the tipping point's length across ordered difficulty quintiles, with right-censored strata assigned maximum rank; the same test repeated on the synthetic tunable-difficulty family's difficulty parameter.

\textbf{Expected outcome.} The trend test rejects the null of no ordering ($p < 0.05$) on both natural and synthetic suites, with the tipping point increasing (or being right-censored) as difficulty increases, and no significant reversal between the two suites.

\textbf{Falsification condition.} Refuted if the trend test finds no significant ordering --- meaning the tipping point is statistically indistinguishable from a single universal length threshold --- or if the ordering reverses between natural and synthetic suites, either of which would mean difficulty-conditioned budget allocation has no basis even if a flip exists somewhere (RQ2).

\textbf{Hypothesis (plain language).} The point where thinking starts to hurt is not a single fixed length for all problems --- it moves with how hard the problem is.

\textbf{Metrics.} Per-stratum tipping-point length with bootstrap confidence interval; trend-test statistic and $p$-value; reversal check between natural and synthetic suites.

\textbf{Result.} \ph

\subsubsection{RQ4: Do cheap trajectory-level signals carry information about the tipping point beyond a length-only null?}
\textbf{Setup.} From the first 256 thinking tokens plus four truncated-answer probes, compute entropy-trajectory and early-consistency features; fit a classifier of whether marginal gain beyond the probe point is negative; compare its AUROC against an identically-trained classifier fed only position-controlled length features, on held-out strata.

\textbf{Expected outcome.} At least one cheap feature family achieves AUROC $\geq 0.65$ (95\% confidence interval excluding 0.5) on held-out strata, and this AUROC is not matched by the length-only null.

\textbf{Falsification condition.} Refuted if no cheap feature family's AUROC confidence interval excludes 0.5, or if the length-only null achieves statistically indistinguishable AUROC --- either would mean the apparent signal is a repackaging of trace length rather than genuine trajectory-level information.

\textbf{Hypothesis (plain language).} What the model's early trace looks like --- not just how much of it exists --- carries real information about whether continuing will help or hurt.

\textbf{Metrics.} Per-feature-family AUROC and its confidence interval; length-only null AUROC for comparison.

\textbf{Result.} \ph

\subsubsection{RQ5: Does the trajectory-level predictor's information about the tipping point exceed what a coarser difficulty-bin decision already recovers?}
\textbf{Setup.} At matched mean token cost, compare accuracy of stopping at the predicted tipping point against the best single fixed budget, a reimplemented DEER-style confidence-exit policy, per-stratum oracle-stop, and a reimplemented three-class difficulty-bin classifier, on held-out strata and held-out task domains (including GPQA-Diamond).

\textbf{Expected outcome.} All three required, matching the pre-registered conjunctive hypothesis: (a) predicted-stop beats the best single fixed budget on paired per-item correctness (McNemar $p < 0.05$); (b) predicted-stop recovers at least 50\% of the oracle-minus-best-fixed-budget accuracy gap at matched cost; (c) the recovered fraction exceeds what the coarser difficulty-bin classifier recovers under the same matched-cost comparison. A partial pass (e.g.\ (b)+(c) without (a)) is reported as such, not rounded up to a full pass.

\textbf{Falsification condition.} Refuted if any of the three conditions fails: no McNemar-significant improvement over the best fixed budget, or no larger oracle-gap share than the coarse difficulty-bin classifier, or (descriptively) no better showing than the DEER-style confidence-exit policy at matched cost --- in which case the trajectory-level signal is redundant with a much cheaper coarse decision or an already-published stopping rule, and the boundary map's practical value is diagnostic (showing where budgets are wasted) rather than actionable as a new stopping policy. This is a claim about the information content of the trajectory-level signal relative to cheaper or already-published alternatives, not a claim that this predictor is the best deployable stopping policy in any absolute sense.

\textbf{Hypothesis (plain language).} Knowing the shape of the early trace adds something a coarse easy/medium/hard label --- or an existing confidence-based stopping rule --- does not already give you.

\textbf{Metrics.} Oracle-gap recovery fraction for predicted-stop vs.\ coarse-bin classifier, at matched mean token cost; McNemar test on paired per-item correctness between predicted-stop and best fixed budget; descriptive Pareto comparison against the DEER-style confidence-exit policy at matched cost.

\textbf{Result.} \ph

\subsubsection{RQ6: Does a detected tipping point replicate across independent length-control manipulations?}
\textbf{Setup.} Among quintiles with a detected tipping point under the truncation arm (RQ2), check whether the forced-continuation arm detects a negative-slope region of the same sign; natural stopping is logged observationally for the same comparison where feasible.

\textbf{Expected outcome.} The negative-slope region replicates in sign under the forced-continuation arm in a majority of quintiles where truncation detects a flip at all; single-arm-only detections are reported separately as unconfirmed.

\textbf{Falsification condition.} Refuted if a majority of truncation-detected flips fail to replicate in sign under forced continuation --- in which case the finding is better explained as an artifact of the truncation-and-forced-answer manipulation than as a property of the reasoning process, and no tipping-point claim from this design can be treated as confirmed.

\textbf{Hypothesis (plain language).} Where thinking starts to hurt is a property of the reasoning process, not just of how the length limit happened to be enforced.

\textbf{Metrics.} Cross-arm sign-agreement rate among quintiles with a detected flip; count of confirmed vs.\ unconfirmed (single-arm-only) detections.

\textbf{Result.} \ph

\subsubsection{RQ7: Does the tipping point coincide with a mechanistic signature of self-overwriting, or only with a generic accuracy decay?}
\textbf{Setup.} Restricted to strata with a detected tipping point under RQ2, compare (i) correct$\to$incorrect vs.\ incorrect$\to$correct answer-flip rates, and (ii) verification/backtracking step density, in the grid-length window just before vs.\ just after $L^*_s$, using a paired bootstrap on the before/after contrast.

\textbf{Expected outcome.} The correct$\to$incorrect flip rate rises significantly after $L^*_s$ relative to before it (paired bootstrap, Holm-corrected across strata, $p<0.05$), and this rise co-occurs with elevated verification/backtracking step density in the same window.

\textbf{Falsification condition.} Refuted if correct$\to$incorrect and incorrect$\to$correct flip rates are statistically indistinguishable before vs.\ after $L^*_s$ --- in which case the localized boundary from RQ2 is real, but its mechanism is not verification-induced self-overwriting; the tipping point's \emph{existence} (RQ2) is unaffected, but the Introduction's causal ``overwrite'' story downgrades to a descriptive co-occurrence (``coincides with,'' not ``is caused by''), and the paper reports the boundary without a mechanism claim.

\textbf{Hypothesis (plain language).} Where thinking starts to hurt is not just a point where accuracy happens to dip --- it is a point where the model's own verification behavior starts overwriting correct answers.

\textbf{Metrics.} Correct$\to$incorrect and incorrect$\to$correct flip rates before vs.\ after $L^*_s$, with bootstrap confidence intervals; verification/backtracking step density before vs.\ after $L^*_s$; correlation between entropy-trajectory slope and flip sign at $L^*_s$.

\textbf{Result.} \ph

\vspace{0.5em}
\noindent\textbf{Further evaluation.} With more time, a natural extension is testing whether the tipping point and its predictor transfer across model families rather than only within the scales and families evaluated here --- explicitly out of scope for this design, since cross-family transfer would require a validation protocol of its own rather than an extension of the held-out-strata/held-out-domain evaluation used above.

\section{Related Work}

The closest prior work to this paper now forms three rings around the same gap, and it is worth naming them before surveying the wider landscape; this ranking was revised in August 2026 after a freshness audit, and the concessions below are correspondingly explicit. Innermost ring --- the finding itself: a systematic study of overthinking in test-time compute scaling \cite{WhenMoreThinkingHurts2604} reports diminishing marginal utility at higher budgets, extended reasoning abandoning previously correct answers, and an optimal thinking length that varies with problem difficulty --- the existence and direction of everything this paper measures, which this design therefore concedes rather than claims. What that study does not provide is a per-stratum boundary statistically localized with stated uncertainty (or right-censoring), evidence that a detected boundary survives independent length-control manipulations, or any pre-hoc predictor of the boundary's location --- precisely the three residual deltas this paper contributes. In the same ring, an inverse-scaling evaluation suite \cite{InverseScalingTTC2507} constructs task categories where longer reasoning deteriorates accuracy and names five distinct failure modes --- existence evidence on constructed tasks, without difficulty stratification on natural benchmarks or boundary localization. Second ring --- shape and pathology: a theoretical account of chain-of-thought length \cite{WhenMoreIsLess6QDF} derives an inverted-U relationship between reasoning length, task difficulty, and accuracy --- predicting exactly the kind of sign flip this paper measures --- but validates that prediction only on synthetic tasks with an experimenter-constructed difficulty variable, never on natural benchmarks under an operational proxy, and offers no per-instance predictor. An empirical study of ``overthinking'' \cite{Chen2024Overthinking} documents that reasoning models expend redundant solution rounds on easy problems with negligible or negative benefit, but reports aggregate efficiency metrics rather than a statistically localized, per-stratum tipping point; an agentic-tasks variant \cite{DangerOverthinking2502} documents the same pathology on a different estimand --- extended internal reasoning displacing environment interaction in interactive software tasks --- rather than within-trace accuracy-versus-length curves; and a survey of efficient reasoning \cite{StopOverthinkingSurvey2503} consolidates this rapidly densifying mitigation literature, none of whose entries publishes a localized boundary. Third ring --- interventions that presuppose the boundary rather than measuring it: representation-steering approaches --- controlling ``thinking speed'' in real time \cite{ThinkingSpeedESEL}, or steering the attention heads implicated in reasoning behavior \cite{251224574} --- concede that an unmodified model's test-time reasoning trajectory is steerable without retraining, and thereby presuppose exactly the difficulty-conditioned boundary at issue, but both report the accuracy and token deltas produced by steering rather than the location or shape of the boundary in the unmodified model, and neither offers a cheap pre-hoc predictor of it. Across all three rings --- existence evidence without localization, a theoretical shape validated only on synthetic difficulty, and interventions that exploit a boundary they never measure --- this paper's contribution is precisely the statistically localized, cross-arm-validated, pre-hoc-predictable boundary measurement none of them provides; its existence claim, by contrast, is conceded to the innermost ring.

\textbf{Confidence-based dynamic stopping.} Training-free early-exit methods \cite{DEERNpU7} detect reasoning-transition boundaries during a single trace and exit once a trial answer's confidence crosses a threshold; distillation-based variants \cite{CaTSjrSc4} fold multi-sample self-consistency agreement into a single-pass confidence predictor for early-stopping repeated whole-trace sampling. Both optimize the decision to stop and validate end-to-end accuracy-cost tradeoffs, but neither publishes the underlying difficulty-conditioned boundary its stopping rule implicitly exploits, and their signals are available only during an ongoing trace (or after a distillation step) rather than as a pre-hoc characterization --- the distillation-based variant additionally operates on a different length axis than this paper (number of resampled generations, not thinking tokens within one trace), so it is cited as evidence of the same unmeasured-boundary gap rather than as a directly comparable stopping-rule baseline.

\textbf{Pre-hoc difficulty and self-consistency probes.} A lightweight probe trained on prefill hidden states \cite{DiffAdapt644F} classifies query difficulty into three coarse regimes before any reasoning token is generated, selecting among fixed reasoning strategies; a related probe \cite{AdaptiveThinking2i6R} instead predicts a query's expected self-consistency as a reasoning-necessity proxy. Both sort queries into pre-hoc regimes cheaply, but neither validates how much of the achievable oracle accuracy-cost tradeoff a coarse decision actually recovers relative to a continuously localized boundary --- precisely the comparison this paper's decoupled cheap-signal predictor is built to report, in either direction.

\textbf{Interventions that presuppose the boundary.} Reinforcement-learning approaches \cite{JET2u5Z} train a model to terminate reasoning once sufficient information has accumulated, using a correctness-conditioned length reward; representation-steering approaches control ``thinking speed'' in real time using a difficulty signal derived from layer-wise logit divergence, either via a contrastive-PCA steering vector combined with real-time difficulty estimation \cite{ThinkingSpeedESEL} or via activation steering of attention heads that encode linear-vs-non-linear reasoning behavior \cite{251224574} --- the latter motivated by its own authors explicitly as a correction for both overthinking and underthinking, the same asymmetry whose boundary this paper measures; and mechanistic studies of overthinking \cite{FirstImpression2PP7} test whether an implicit early-answer bias causally triggers excessive reflection. Each of these changes the model or its decoding to exploit an assumed difficulty-conditioned boundary, entangling the intervention with the phenomenon it presupposes --- none measures where the boundary lies in an unmodified model, which is the measurement this paper isolates before any intervention is considered. Both steering-based interventions \cite{ThinkingSpeedESEL,251224574} in particular presuppose a difficulty-conditioned control signal without independently validating that the signal predicts a statistically localized tipping point, which is exactly the validation this paper's decoupled cheap-signal predictor supplies.

\textbf{Budget forcing and uniform length control.} Externally truncating a thinking segment or appending a continuation cue to force further reasoning \cite{Muennighoff2025s1} demonstrates that thinking length is a controllable scaling knob, with reported scaling curves (and their flattening) taken in aggregate. Applied uniformly across queries, this approach cannot say for which problems the forced extra thinking is the part that hurts --- the difficulty-conditioning question this paper's stratified design is built to answer, and the uniform-budget-forcing arm serves directly as this paper's no-stratification control. An aggregate compute-accuracy Pareto analysis across open-source reasoning models \cite{251224776} independently documents a rising-then-saturating curve at the population level, without difficulty stratification or a localized sign flip --- the same no-stratification limitation this paper's fixed-budget baseline is designed to make concrete at the per-stratum level.

\section{Conclusion}

This work reframes a budget-allocation assumption --- that more test-time thinking is safe to keep adding until it simply stops helping --- as a boundary-measurement question with a precise design, building on recent evidence that the failure mode exists \cite{WhenMoreThinkingHurts2604,InverseScalingTTC2507} and conceding that existence to prior work rather than re-claiming it, while supplying what that evidence still lacks: localization with stated uncertainty, manipulation-artifact control, and a pre-hoc predictor. The central methodological wager is that a tipping point is only trustworthy once the difficulty stratification it's conditioned on is cross-validated, the change point is localized with paired statistical power rather than eyeballed, and the finding survives at least two independent ways of imposing a length limit --- weaker designs risk reporting a manipulation artifact as a property of reasoning itself. Whichever regime the data show, the finding is decision-relevant: a detected, difficulty-conditioned tipping point would mean uniform or ever-growing thinking budgets are actively harmful on a predictable subset of problems, while a well-powered null within the studied budget range would itself bound where such harm cannot occur, refuting the premise several existing stopping-rule and length-penalty methods presuppose. Two limitations bound this claim from the outset and are natural next steps rather than afterthoughts: no experimental design here fully identifies the causal counterfactual of an unmodified model choosing to think longer, so any causal reading of the tipping point remains qualified by cross-arm agreement rather than proven by it; and the cheap-signal predictor's cross-model-family transfer is explicitly out of scope, leaving open whether the trajectory-level information found within one model family generalizes to reasoning models built differently.

\vspace{1em}
\hrule
\vspace{1em}

\section*{Pre-handoff checklist}
\begin{itemize}[leftmargin=1.4em]
\item \textbf{Organization/navigation} --- narrative flows problem $\to$ method $\to$ experiments $\to$ related work without forward references to undefined terms.
\item \textbf{Terminology discipline} --- ``difficulty stratum,'' ``tipping point,'' ``right-censored,'' ``oracle-gap recovery'' each defined at first use in Problem Definition, used consistently thereafter; no internal stage IDs (\texttt{sp\_*}, \texttt{m\_*}, \texttt{rq\_*}, \texttt{C1}-\texttt{C4}, \texttt{M1}-\texttt{M5}) leaked into prose --- checked directly against the discussion pipeline log.
\item \textbf{Notation with worked example} --- the formal setup section defines $(f, x, y, s(x), L, \mathcal{L}, Z, A_s(L), L^*_s)$ and the Problem Definition/Methodology sections consistently reuse this notation; no separate worked numeric example is included since no numbers exist yet (would require fabricating results).
\item \textbf{RQ framing matches outcomes} --- all 7 RQs are graded/quantitative (correlations, rates, AUROCs, recovery fractions, flip-rate contrasts), each answered in its own subsection with an explicit falsifiable expected outcome and an explicit falsification condition, not a binary yes/no. RQ7 (mechanism) closes a gap surfaced at Phase G review: paper.md's M6 mechanism analysis (guidance.md Part 2 A2's $\geq$1-mechanism-analysis evidence standard) had no corresponding RQ in an earlier draft of this file.
\item \textbf{No typesetting artifacts} --- only explicit placeholder markers in Results, per the drafting convention this is a pre-registered spec, not a report.
\item \textbf{Artifact plan} --- the length-controlled generation harness with per-token logprobs, the change-point analysis code, and the trained cheap-signal predictor are the natural release artifacts; not yet built, so not claimed as available.
\end{itemize}

\section*{Hand-off}
This is a draft, not a finished paper --- experiments have not run. Natural next steps: prose polish and venue-specific LaTeX formatting once a target venue is chosen, or revising an earlier stage if piloting surfaces a problem (e.g., if RQ1's proxy-agreement rate is much lower than expected).

\bibliographystyle{plain}
\bibliography{refs}

\end{document}
